% The TLS heartbeat request/response (RFC 6520) as a byte-field strip.
% Used in M10-L05.
\documentclass[tikz,border=4pt]{standalone}
\usepackage{tikz}
\input{_m10-styles.texinput}
\begin{document}
\begin{tikzpicture}[m10, x=1mm, y=-1mm, font=\large]
  % request
  \node[note, anchor=east, font=\Large\bfseries] at (-3,8) {request};
  \node[stackcell, minimum width=20mm, minimum height=12mm, font=\large] (t)   at (10,8)  {type};
  \node[stackcell, minimum width=28mm, minimum height=12mm, font=\large] (len) at (34,8)  {length};
  \node[heapcell,  minimum width=42mm, minimum height=12mm, font=\large] (pay) at (69,8)  {payload};
  \node[stackcell, minimum width=28mm, minimum height=12mm, font=\large] (pad) at (104,8) {padding};
  \node[note, anchor=north, font=\large] at (10,15)  {1 byte};
  \node[note, anchor=north, font=\large] at (34,15)  {2 bytes};
  \node[note, anchor=north, font=\large] at (69,15)  {length bytes};
  \node[note, anchor=north, font=\large] at (104,15) {\(\geq 16\)};

  % response echoes the payload
  \node[note, anchor=east, font=\Large\bfseries] at (-3,38) {response};
  \node[heapcell, minimum width=42mm, minimum height=12mm, font=\large] (rpay) at (69,38)
    {same payload, echoed};
  \draw[goodflow, line width=1pt] (pay.south) -- (rpay.north);

  \node[note, anchor=north west, font=\large\itshape] at (-20,48)
    {``Send back these \texttt{length} bytes.'' The server must trust only what it received.};
\end{tikzpicture}
\end{document}
